\documentclass[11pt]{article}
\usepackage[margin=1in]{geometry}
\usepackage[T1]{fontenc}
\usepackage{lmodern}
\usepackage{amsmath,amssymb,amsthm}
\usepackage{graphicx}
\usepackage[hidelinks]{hyperref}
\usepackage{microtype}
\emergencystretch=2em

\newtheorem{theorem}{Theorem}
\newtheorem{lemma}{Lemma}
\newtheorem{remark}{Remark}

\newcommand{\A}{\mathcal A}
\newcommand{\B}{\mathcal B}
\newcommand{\N}{\mathbb N}
\newcommand{\eps}{\varepsilon}
\newcommand{\op}{\operatorname{op}}
\newcommand{\Id}{\operatorname{Id}}

\title{A \(B(c_0)\) coefficient separating left and right Asplund/tame functionals}
\author{Run fa\_banach\_001, agent\_super\_01}
\date{2026-06-30}

\begin{document}
\maketitle

\section*{Source Problem}

The source paper is Matan Komisarchik,
\emph{Characterizations of Asplund and Tame Functionals using Arens Products},
arXiv:2602.14317.  In Definition 2.29 it defines, for a Banach algebra
\(\A\) and \(\varphi\in \A^*\),
\[
 \varphi\in \operatorname{Asp}^{(l)}(\A)
     \quad\Longleftrightarrow\quad
 B_\A\circledast\varphi\text{ is an Asplund subset of }\A^*,
\]
\[
 \varphi\in \operatorname{Asp}^{(r)}(\A)
     \quad\Longleftrightarrow\quad
 \varphi\circledast B_\A\text{ is an Asplund subset of }\A^*,
\]
and analogously for left and right tame functionals.  Remark 2.30 states that
the relationship between left and right Asplund or tame functionals is still
open.

\begin{center}
\includegraphics[width=0.98\linewidth]{figures/open_problem_crop.png}
\end{center}

\section*{Answer}

\begin{theorem}\label{thm:main}
There is a unital Banach algebra \(\A\) and a functional
\(\varphi\in\A^*\) such that \(\varphi\) is left Asplund and left tame, but
\(\varphi\) is neither right Asplund nor right tame.  Consequently the left
and right notions are not equivalent.  Passing to the opposite Banach algebra
also gives a functional which is right Asplund and right tame but not left
Asplund or left tame, so neither implication holds in general.
\end{theorem}

\section*{Construction}

Let \(X=c_0\), let \(\A=\mathcal L(X)\) with composition, and write \(e_1\)
for the first unit vector of \(c_0\).  Let \(\epsilon_1=e_1^*\in X^*=\ell_1\)
be the first coordinate functional.  Define the matrix coefficient
\[
   \varphi(T)=\epsilon_1(T e_1), \qquad T\in \mathcal L(c_0).
\]
The source convention for the dual bimodule is
\[
  \langle S\circledast\varphi,T\rangle=\langle\varphi,TS\rangle,
  \qquad
  \langle \varphi\circledast S,T\rangle=\langle\varphi,ST\rangle .
\]

For \(y\in X\) and \(\eta\in X^*\), define
\[
   J_X(y)(T)=\epsilon_1(Ty),
   \qquad
   J_{X^*}(\eta)(T)=\eta(Te_1).
\]

\begin{lemma}\label{lem:orbits}
The two orbit sets are exactly
\[
   B_\A\circledast\varphi=J_X(B_X),
   \qquad
   \varphi\circledast B_\A=J_{X^*}(B_{X^*}).
\]
\end{lemma}

\begin{proof}
For \(S\in B_\A\) and \(T\in\A\),
\[
  (S\circledast\varphi)(T)=\varphi(TS)=\epsilon_1(T(S e_1))
       =J_X(S e_1)(T).
\]
Since \(S e_1\in B_X\), this proves one inclusion.  Conversely, if
\(y\in B_X\), then
\[
   S_y z=\epsilon_1(z)y
\]
defines an operator \(S_y\in B_\A\), and \(S_y e_1=y\).  Hence
\(J_X(y)=S_y\circledast\varphi\).

Similarly,
\[
  (\varphi\circledast S)(T)=\varphi(ST)=\epsilon_1(S(Te_1))
       =(S^*\epsilon_1)(Te_1)=J_{X^*}(S^*\epsilon_1)(T).
\]
For the reverse inclusion, given \(\eta\in B_{X^*}\), define
\[
   R_\eta z=\eta(z)e_1 .
\]
Then \(R_\eta\in B_\A\) and \(R_\eta^*\epsilon_1=\eta\).  Thus
\(J_{X^*}(\eta)=\varphi\circledast R_\eta\).
\end{proof}

\begin{lemma}\label{lem:isometric}
The maps \(J_X:X\to\A^*\) and \(J_{X^*}:X^*\to\A^*\) are isometries.
\end{lemma}

\begin{proof}
For \(y\in X\), the upper bound is immediate:
\[
  |J_X(y)(T)|=|\epsilon_1(Ty)|\le \|T\|\|y\|
       \qquad (T\in B_\A).
\]
For the reverse inequality, choose \(f\in B_{X^*}\) with
\(|f(y)|>\|y\|-\eps\), and set \(T_f z=f(z)e_1\).  Then \(T_f\in B_\A\) and
\[
   |J_X(y)(T_f)|=|\epsilon_1(f(y)e_1)|=|f(y)|>\|y\|-\eps .
\]
Letting \(\eps\downarrow0\) gives \(\|J_X(y)\|=\|y\|\).

For \(\eta\in X^*\), again
\[
  |J_{X^*}(\eta)(T)|=|\eta(Te_1)|\le \|\eta\|\|T\|
       \qquad (T\in B_\A).
\]
Choose \(v\in B_X\) with \(|\eta(v)|>\|\eta\|-\eps\), and define
\(T_v z=\epsilon_1(z)v\).  Then \(T_v\in B_\A\), \(T_v e_1=v\), and
\[
   |J_{X^*}(\eta)(T_v)|=|\eta(v)|>\|\eta\|-\eps .
\]
Thus \(J_{X^*}\) is also an isometry.
\end{proof}

\section*{Asplund and Tame Status of the Two Orbits}

The source paper recalls the standard facts that a Banach space \(V\) is
Asplund if and only if \(B_V\) is an Asplund subset, and \(V\) is Rosenthal
if and only if \(B_V\) is tame.  It also uses that these subset classes are
preserved by bounded linear maps; in particular they are invariant under
linear isomorphism onto their closed spans.

The space \(c_0\) is Asplund, since \(c_0^*=\ell_1\) is separable.  It is also
Rosenthal: every subspace of \(c_0\) has separable dual, while \(\ell_1^*
=\ell_\infty\) is nonseparable, so \(c_0\) cannot contain a copy of
\(\ell_1\).  Hence \(B_{c_0}\) is both Asplund and tame.  By
Lemmas~\ref{lem:orbits} and~\ref{lem:isometric},
\[
   B_\A\circledast\varphi=J_X(B_X)
\]
is an isometric copy of \(B_{c_0}\), so it is Asplund and tame.  Therefore
\(\varphi\) is left Asplund and left tame.

On the other hand \(X^*=\ell_1\) is not Asplund, because it is separable but
its dual \(\ell_\infty\) is nonseparable.  It is not Rosenthal, because it
contains a copy of \(\ell_1\), namely itself.  Hence \(B_{\ell_1}\) is neither
Asplund nor tame.  Since
\[
   \varphi\circledast B_\A=J_{X^*}(B_{X^*})
\]
is an isometric copy of \(B_{\ell_1}\), this orbit is neither Asplund nor
tame.  Therefore \(\varphi\) is not right Asplund and not right tame.

Finally, let \(\A^{\op}\) be the opposite Banach algebra.  The same Banach
space and functional have left and right orbit sets interchanged in
\(\A^{\op}\).  Thus the opposite algebra gives the reverse separation:
right Asplund and right tame need not imply left Asplund or left tame.
\(\square\)

\section*{Verification Notes}

The proof uses no numerical computation.  The accompanying script
\texttt{code/check\_orbit\_formula.py} checks the two orbit identities in
finite-dimensional matrix algebras, which is only a sanity check for the
algebraic formulas in Lemma~\ref{lem:orbits}.

The source paper notes immediately after the open-question remark that it
uses all functionals in \(\A^*\), unlike a convention in the earlier
tame-functional paper where one restricts to left or right uniformly
continuous functionals.  This example uses the source definitions.  Since
\(\mathcal L(c_0)\) is unital, the usual uniformly continuous restriction
should not remove the example, but this convention-specific point is marked
for human review.

\section*{Novelty Check}

Before promotion, the run's cheap indexes were searched for arXiv:2602.14317,
the exact source title, the phrases ``left and right Asplund'', ``left and
right tame'', ``tame functionals'', and \(B(c_0)\)-based variants.  No prior
packet or attempt for this counterexample was found.  Local source-corpus
searches found the source Remark 2.30 and the later Example 5.5, which says
that for \(V=c_0\) and \(\A=\mathcal L(V)\), all matrix coefficients are left
Asplund.  That source example does not identify the right orbit or give the
separating counterexample above.

Web spot checks on 2026-06-30 for the exact source title, arXiv:2602.14317,
``Quest. 7.2'' with ``tame functionals'', ``left right Asplund tame
functionals'', and ``\(B(c_0)\) tame functionals Asplund'' found the source
problem and related tame-functional literature, but no exact later resolution
or this \(B(c_0)\) coefficient construction.

\section*{Human Review Recommendation}

The recommended review focus is narrow: verify the source-paper conventions
for Asplund/tame subsets and the standard invariance of these subset classes
under isometric embeddings.  Once that is accepted, the two orbit calculations
show the separation directly.

\section*{References}

\begin{itemize}
\item Matan Komisarchik, \emph{Characterizations of Asplund and Tame
Functionals using Arens Products}, arXiv:2602.14317, 2026.
\item Matan Komisarchik, \emph{The Banach algebra \(L^1(G)\) and tame
functionals}, Comment. Math. Univ. Carolin. 65(2), 131--158, 2024.
\item E. Glasner and M. Megrelishvili, \emph{Representations of dynamical
systems on Banach spaces}, and related Rosenthal/tame-family results cited as
Fact 2.23 in the source paper.
\end{itemize}

\end{document}
